\chapter[Judaism: Covenant, Law, and Diaspora]{Judaism: Covenant, Law, and Community in Diaspora}
\section{A Small Box on a Doorpost}
Begin with a small object. In Jerusalem's Old City, Brooklyn, north London, Buenos Aires apartment buildings, wherever Jewish people live, you may find it on a doorway's right side: a finger-length box of wood, metal, or ceramic, fixed at an angle. Some entering or leaving touch it and kiss their fingertips.

Called a mezuzah, it contains no jewels or charm, only a small parchment scroll with two Hebrew passages carefully written by a trained scribe. Its opening says:

\begin{quote}
``Hear, Israel! The Lord our God is the one Lord.'' (Deuteronomy 6:4)
\end{quote}
It continues: love this one God with heart, soul, strength; keep these words within, teach children diligently, speak them at home, on roads, lying down, rising; bind them on hand and forehead, write them on doorposts and gates (Deuteronomy 6:5--9).

The small box follows that instruction to write on doorposts.

Notice its distinctiveness. Most peoples locate identity in land, ancestry, borders, capitals. Here identity is copied onto parchment, nailed to doors, bound to arms, recited, carried everywhere. Land may fall, gates burn, kings die, but a scroll can travel under an arm.

Judaism is among humanity's oldest monotheistic faiths, affirming one God. Yet it is more: book, law, contract, weekly calendar, community persisting for almost two millennia without a national territory. How did this work? The story begins with fire.

\section[70 CE: A Coffin through the Gates]{70 CE: The Coffin Carried out through the City Gates}
Come into a scene.

Summer, 70 CE, Jerusalem besieged for months by Roman troops under imperial son Titus. Food has run out. Witness-historian Josephus describes people eating belts, roots, worse. Smoke and death fill the air. Defenders fight Rome and one another. Radical Zealots guard the gates, preferring collective starvation to surrender, which betrays God.

For centuries the temple concentrated Jewish life's greatest weight: the one God's earthly dwelling, three yearly national pilgrimages, priestly sacrifice, continual incense. Imagine the people's sole heart, synagogues only branches. Rome now encircles it.

An elderly teacher, Yohanan ben Zakkai, acts in a story told by later rabbinic literature, the \textit{Talmud}, ``Gifts,'' 56b. Written centuries afterwards and likely legendary in details, it precisely expresses rabbinic self-understanding:

Yohanan judges city and temple indefensible; fighting Rome will bury the people with them. Pupils place him in a coffin carried out under burial rules. At the Roman camp, the dead man sits up and asks for the commander. He tells Vespasian, Titus's father, that he will become emperor. Soon news reports Nero dead and armies acclaiming Vespasian, who prepares to return for accession. Offered a reward, Yohanan asks neither city pardon nor temple preservation but Yavneh and its scholars.

``Give me Yavneh and its sages,'' the Talmud says in essence.

Rome agrees. Late that summer the temple burns and great stones collapse. Josephus describes daylight-bright night and cries louder than flames. It is never rebuilt. In coastal Yavneh, teachers begin something apparently useless: word-by-word discussion of how to live this faith without its temple.

\section{Without the Temple, What Remains?}
State the problem before answering.

What remains when religion loses temple, high priest, land, king? Ancient precedent often says little. Empires conquer, deport elites, remove images; generations later names and languages dissolve like salt into conquerors' soup. After Assyria deported northern Israelite elites, the ten tribes largely disappear from records, becoming the lost tribes. Not stupidity, but a small people's common imperial fate.

After 70 CE, Jewish history refused that script. Why?

First recover origins: who were these people, and where did their book come from?

Hebrew biblical narrative begins with Abraham, a herdsman called from home into covenant: an enduring reciprocal contract promising descendants and land in return for belonging and obedience to God. Centuries later, descendants enslaved in Egypt leave under Moses, receive Sinai law, renew covenant. Later Canaanite kingship produces David and Solomon, who builds Jerusalem's first temple.

Honesty requires that this early account survives almost entirely in the Hebrew Bible. Archaeology has found no direct evidence of Exodus or Sinai covenant; scholars long dispute occurrence, scale, correspondence with the narrative. \textbf{Traditional belief differs from historical evidence}. Tradition affirms Exodus and reenacts it annually; evidence establishes its central place in communal memory over two millennia ago. For understanding religion, that may matter more.

Evidence becomes firmer after a point. Following the tenth century BCE, the kingdom divides. Northern Israel falls to Assyria in 722; southern Judah lasts over a century more before Neo-Babylonian conquest. In 586, Nebuchadnezzar's forces burn the first temple and deport royalty, priests, craftsmen. Cyrus takes Babylon in 539. His conciliatory policy, reflected in the Cyrus Cylinder's returns and temple restorations, benefits Judeans among others. Some return and rebuild in the late sixth century.

More important follows. Nehemiah 8 describes Ezra reading Moses's law at the Water Gate from dawn until noon, with explanation enabling understanding. People weep, many hearing the complete provisions for the first time. Scholars generally locate collection, editing, and stabilisation of biblical cores, especially the Five Books or Torah, teaching, in exile and return. Court and temple archives become what all must hear and observe.

Notice the change: the community's constitution leaves palace locks for weekly public reading. Kings die, temples burn, reading needs no king's permission.

Several enduring elements develop. \textbf{Monotheism}: one God, others not merely foreign gods but no gods, unusual amid polytheistic Egypt, Babylon, Greece. \textbf{Covenant}: blessings for observance, penalties for breach; temple destruction means broken commitments, not divine defeat, making failure fuel faith. \textbf{Law}: traditionally 613 Torah commandments encompass murder, food, clothing, rest, all daily life. \textbf{Sabbath}: seventh-day cessation from Friday sunset to Saturday sunset, as Exodus 20 commands remembrance and sanctification. Amid production's universal bustle, a people cuts out weekly time for rest, reunion, God.

This is Yohanan's whole portable inheritance, carried into and out of the coffin. Could it replace a temple, or merely choose a slower collective death?

\section{Two Ways to Live}
The destroyed-temple question was already a second examination in 70 CE. The first, 586 BCE, left two opposed biblical voices.

\textbf{First: Remember and never assimilate.} Psalm 137 speaks from Babylonian exile:

\begin{quote}
``By Babylon's rivers we sat and wept, remembering Zion. We hung our harps upon the willows \,\ldots\, How can we sing the Lord's song in a foreign land?''
\end{quote}
Zion is Jerusalem. Captors seek entertaining songs; exiles hang instruments away. Their song is unmarketable homesickness, refusing absorption, forgetting, comfortable enemy-land living.

\textbf{Second: Settle and live well.} Around the same time, Jeremiah writes from Jerusalem to Babylonian exiles in Chapter 29: build and inhabit houses, plant and eat, marry and raise children, seek that city's peace, because its peace brings theirs.

Weigh the letter. Its author will suffer the collapse and hates Babylon deeply, yet urges not immediate-homecoming obsession but life-building and prayer for an enemy city. Not surrender, but longer loyalty to survival rather than present posture.

The voices alternate for two millennia. In 70 they argue through Zealots and Yohanan.

Zealots and later Bar Kokhba followers, 132--135, easily look recklessly suicidal today. Fairness requires their full case. Temple is divine cohabitation, not mere stone; abandoning it betrays covenant and leaves merely biological living. Surrender guarantees no survival, only slower assimilation, visible across Roman peoples. And postponed freedom often never arrives. None is foolish. The two great revolts cost Rome several legions. They lost, and history lacks patience with losers.

Yohanan's side likewise argues fully: lives outrank walls, books weigh less than stone and travel. While people read, teach, discuss, temple survives relocated. History gives outcomes, not grades: Masada's fortress becomes a memorial, Yavneh's school the birthplace of two millennia of Judaism.

Another voice matters. Disputes concern faith itself, not only survival abroad. While the temple stood, Amos spoke for God almost rebelliously:

\begin{quote}
``I hate your festivals and take no pleasure in your solemn assemblies.'' (Amos 5:21)
\end{quote}
Micah states it more clearly:

\begin{quote}
``Humanity, the Lord has shown what is good. What does he require? Act justly, love mercy, walk humbly with your God.'' (Micah 6:8)
\end{quote}
Priests locate faith in temple rites and smoke; prophets say injustice makes perfect ritual stink. Internal protest, not foreign attack. A tradition canonising and repeatedly reading its critics is no seamless stone.

\section[Thinking Bridge: Four Layers of Reading]{Thinking Bridge: Read Four Layers in One Scriptural Page}
After 70, rabbis, my teachers, become a new authority replacing priests. Sacrifice is unavailable, Torah remains. They relocate religion into text: read sacrificial passages, study in synagogues. Learning becomes worship. Their portable reading practice later receives four levels, initials spelling PaRDeS, orchard:

\begin{itemize}
\item \textbf{Literal (Peshat)}: what the text plainly says.
\item \textbf{Allusive (Remez)}: echoes and hints; why this word rather than its synonym, and which other passage it gestures towards.
\item \textbf{Interpretive (Derash)}: inference, extension, filling gaps; what to do in unmentioned cases.
\item \textbf{Secret (Sod)}: hidden mysteries, chiefly explored by later Kabbalah.
\end{itemize}
Not a rigid workflow but a habit: \textbf{important texts deserve four readings with different questions}. Consider Torah's famous provision:

\begin{quote}
``Eye for eye, tooth for tooth.'' (Exodus 21)
\end{quote}
Literally, blinding another costs an eye. Before calling it cruel, remember blood feuds where one injury killed a clan. This was an \textbf{upper limit}, retaliation no greater than injury, brake rather than accelerator.

Rabbis pursue interpretation. What if the offender has one eye and losing it makes him wholly blind? How account for severe versus slight injury? How ensure exact equality in extraction? Talmudic arguments conclude \textbf{monetary compensation}, assessing lost work, treatment, pain, humiliation. Eye for eye becomes a compensation schedule.

The reading never declares sacred words obsolete or deletes them, but allows meanings to be debated so extensively that literal sense is transformed. Unchanged wording, changed living. Critics call it sophistry; rabbinic self-understanding in an earliest Mishnah chapter, Avot 1:1, instructs a fence round Torah, a low outer barrier preventing crossing the real boundary.

Use it today. No snacks at school literally regulates mouths, but interpretation asks about hunger between lessons, birthday cakes, intended harms. Laws, contracts, even I'm fine invite four questions. Shallow readers ask what it says; deeper readers add what it echoes, what its silence requires, what lies beneath.

\section[Reading the Words on the Doorpost]{Reading Evidence: The Words on the Doorpost}
Read the opening box's text, the Shema, Hebrew for hear, recited morning and evening:

\begin{quote}
``Hear, Israel! The Lord our God is the one Lord. Love the Lord your God with all heart, soul, and strength. Keep today's commanded words in your heart; teach them diligently to your children; speak of them at home and on the road, lying down and rising.'' (Deuteronomy 6:4--7)
\end{quote}
Notice three things in this ancient text.

First, it begins hear, not believe. Hearing is daily action. The tradition centres not agreement with propositions alone but what is done daily.

Second, memory is stored in hearts, children's ears, homes, roads, lying and rising, every time. Weekly religious visits do not suffice; faith occupies daily gaps. Many of the 613 commandments turn food, clothing, weekly routines into reminders. Identity is rehearsed at every meal.

Third, children are named as education's recipients. The next generation is not automatically us; transmission is daily labour. Literacy later becomes prominent because reading law is commanded and fathers teach sons Torah. A farming and artisan people making literacy religious duty was unusual, later a great resource.

Look back at the mezuzah: no decoration, but an outside terminal of continuous self-reminding.

\section[Why Did This Community Survive?]{Why Did They Remain When So Many Disappeared?}
At least three explanations address vanished ancient peoples versus persistent Jewish communities, each partial.

\textbf{First: covenant, the tradition's account.} God never breaks agreement. Captivity, exile, burned temple discipline breaches, not abolish covenant. Repentance renews observance. Each disaster gains meaning, not defeated deity but outstanding debt, strengthening rather than dissolving faith. Its limit: internal explanation untestable outside, unable to explain why equally devout peoples did not similarly survive.

\textbf{Second: boundaries, sociology's account.} Survival amid others depends on clear enforceable borders, not inward ideas alone. Sabbath marks time, food laws the table, circumcision the body, marriage kinship. Add literacy and a minyan of any ten adult males, a low organisational threshold, and the result is \textbf{a transplantable community}. Ten households can reproduce synagogue, school, Sabbath, festivals anywhere. Sharp and testable, but reducing faith to survival technique cannot explain dying for boundaries. Without covenant's meaning, people find fences obstructive and leave.

\textbf{Third: surroundings, history's account.} Environment bears half the account. Cyrus's return policy reflects Persian conciliation, not only Jewish resilience. Medieval Islamic societies broadly permitted communities, flourishing in al-Andalus and Ottoman Salonica, while Christian Europe repeatedly expelled them: England 1290, France 1394, Spain 1492. Ottomans received Spanish refugees, Sephardim rooting in Salonica and Istanbul. Survival does not arise wholly internally; slightly looser or tighter surroundings alter outcomes. Its limit is environmental determinism: the same wind overturns some boats, not others, so inspect the boat too.

A corrective to inevitable miraculous survival is \textbf{Kaifeng Jewry}. From around Song times, merchants settled there, building a synagogue in 1163, reading Hebrew, keeping Sabbath, circumcision, food rules. All boundary devices existed. Yet society offered no religious hostility and routes upward through study, examinations, office. Over centuries rabbis died without successors, Hebrew faded, the synagogue deteriorated. By the nineteenth century the community largely assimilated, leaving inscriptions and descendants' faint memory.

Kaifeng qualifies every explanation. Complete fences can still disappear when a welcoming mainstream does not exclude; insiders dismantle them. Survival is no destiny but choices under particular circumstances. Less mythic, more historically precarious.

One further misconception: encountering Judaism through the Christian Old Testament can make it seem an unfinished prequel awaiting completion. That views it from another household's door. After Christianity emerged in the first century, Judaism lived creatively for two more millennia: Mishnah around 200, Babylonian Talmud in fifth--sixth centuries, Maimonides's twelfth-century legal synthesis, Spanish medieval poetry and philosophy, eighteenth-century Hasidism. Still living and arguing, it is poorly served by an Old Testament label, like calling an active old author merely a newcomer's predecessor and missing his whole later career.

\section[Further Challenge: A People in a Book]{Further Challenge: Can a Book Contain a People?}
Now introduce a theory and challenge it.

Benedict Anderson's \textit{Imagined Communities} (1983) calls a nation an imagined political community. No Chinese person knows the other billion-plus, yet all share a picture of belonging. Print matters: common-language newspapers and novels give strangers simultaneity, bringing nationhood into being.

Test Judaism against it. Anderson emphasises printing, but Jewish versions predate Gutenberg by over a millennium using manuscripts. Their substitute is \textbf{synchronised recitation}: weekly Torah readings through an annual cycle, the same passage on the same Sabbath in Krakow, Baghdad, Marrakesh, prayers oriented towards Jerusalem. Geographically fragmented, temporally aligned. Heinrich Heine's frequently quoted image makes the Bible a portable homeland. A homeland in pages cannot be lost.

This adds rather than refutes: \textbf{manuscripts and fixed recitation rhythms can perform the same work as print capitalism}. An older case improving theory is an honour; good theories welcome reality's amendments.

Durkheim adds that worship repeatedly affirms collective existence. Passover dinners exemplify it: fixed procedures, bitter herbs recalling slavery, unleavened bread hurried departure, youngest children's questions about this different night, collective Exodus telling. Each must imagine \textbf{personally} leaving Egypt. Knowing ancestral history is insufficient; yearly reenactment turns collective memory personal. Twenty years of dinners embeds us bodily without theological expertise.

Boundaries are therefore fences and glue, separating outside while repeatedly joining inside. From one side wall, from another roof. Carry both views to the ending.

\section[Switching Off and Asking Who We Are]{Today: A Day Switched Off, and New Arguments about Identity}
This ancient tradition lives in your era, still arguing internally.

Today Jews live in Israel and worldwide through many forms. Orthodox observance strictly follows the 613 commandments, avoiding electricity and travel on Sabbath. Reform, arising in nineteenth-century Germany, adapts law's spirit, seats women and men together, permits women rabbis. Conservative Judaism stands between. Many secular Jews avoid worship yet identify, celebrate Passover, tell Jewish jokes, support schools. Israel's 1948 founding restored statehood after nearly nineteen centuries of diaspora without ending disputes. The Law of Return grants Jewish immigration, but who is Jewish: a Jewish mother's child, a convert, an ethnic rather than religious identifier? Ancient covenant still argues with registration law.

Diversity also includes Ashkenazim of central and eastern Europe, Spanish Sephardim, Middle Eastern and North African Mizrahim, Ethiopian Beta Israel, Kaifeng descendants. Language, food, appearance, rites vary. No pure bloodline binds them; book, calendar, and generations of unfinished questions do.

A gift to everyone is Sabbath. Friday sunset to Saturday sunset suspends work and trade for meals, reading, walks, a weekly collective shutdown in a nonstop world. Psychologists and programmers advocate digital Sabbaths, one phone-free day. Anyone trying knows the hand's unconscious pocketward movement. A three-thousand-year-old shutdown feels harder than mountain-climbing to algorithm-raised people, showing its grasp of a fundamental weakness. Production never ends; humans need compulsory permission to rest. The command removes even the right to overtime, ancient labour protection still relevant.

Resize for yourself. Without ancient faith, you still belong to class, team, dialect community, family table. A new school invites choosing old-friend chats like Psalm 137 or Jeremiah's planting and civic peace. Keep dialect only with grandparents or abandon it? Give a group insider codes? Each votes between remembering and settling. Babylon's debate reopens daily in class.

\section{Your Turn: Boundary or Bridge?}
Return to Kaifeng's synagogue and European ghettos.

Kaifeng's open welcome ends over generations with vanished synagogue, unread texts, communal absorption: nobody harmed, a community gone. European external walls preserve Sabbath, synagogue, scroll, us, but block escape as well as cultural loss. Expulsion and murder recur, culminating in the twentieth century.

Do boundaries protect or imprison? Sabbath and food law preserve a people, fences succeeding, yet identify who can be expelled as not belonging. Kaifeng shows assimilation can require no malice, only welcome and time.

Consider any group of yours, family, class, enthusiasts, people. Its us/them line may be rule, accent, threshold. When does it protect, when imprison? Who should hold the criterion, insiders or outsiders?

Answer with reasons, then face one last question: if your answer changes when you cross the line, which version of you gets to decide?
